\chapter{Vaginal Irritation and Infection}

\section*{Definition}
Inflammation of the vaginal mucosa characterized by itching, burning, abnormal discharge, odor, and discomfort, most commonly caused by infectious or irritant triggers.

\section*{Pathophysiology}
Vaginal ecosystem disruption alters normal lactobacillus-dominated flora, increasing pH and allowing pathogen overgrowth. Common pathogens: Candida albicans (fungal), Gardnerella vaginalis with other anaerobes (bacterial vaginosis), Trichomonas vaginalis (parasitic). Non-infectious causes include chemical irritants, allergic reactions, and atrophic changes.

\section*{Causes}
\begin{itemize}
  \item Vulvovaginal candidiasis (yeast infection)
  \item Bacterial vaginosis
  \item Trichomoniasis (sexually transmitted)
  \item Chemical irritants (soaps, douches, detergents, spermicides)
  \item Atrophic vaginitis (menopause-related)
  \item Allergic reactions (latex, lubricants, contraceptive devices)
  \item Poor hygiene or excessive hygiene
  \item Tight non-breathable clothing
  \item Diabetes mellitus (increased glucose promotes yeast)
  \item Antibiotic use disrupting normal flora
\end{itemize}

\section*{When to See a Doctor}
See your doctor if:
\begin{itemize}
  \item Severe or worsening symptoms
  \item Vaginal irritation with abnormal discharge, fever, pelvic pain, or recurrent episodes
  \item Accompanied by other concerning signs
\end{itemize}

\section*{Related Symptoms}
\begin{itemize}
  \item Vaginal discharge
  \item Itching
  \item painful urination
  \item Painful sex
  \item Abnormal odor
\end{itemize}
\textbf{Important:} If this symptom persists, your doctor will check for serious underlying causes.

\section*{Self-Care / Home Management}
\begin{itemize}
  \item Rest and avoid activities that make it worse.
  \item Maintain adequate hydration and a balanced diet.
  \item Over-the-counter remedies may provide symptomatic relief (consult a pharmacist or doctor).
  \item Monitor symptoms and seek professional advice if they persist or worsen.
  \item Avoid self-medicating with prescription medications without medical guidance.
\end{itemize}

\section*{How Your Doctor Will Evaluate You}
\begin{itemize}
  \item A thorough history and physical examination are the first steps in evaluation.
  \item Laboratory tests (blood work, urinalysis) may be ordered based on clinical suspicion.
  \item Imaging studies (X-ray, ultrasound, CT, MRI) may be indicated when structural pathology is suspected.
  \item Specialist referral is arranged when the diagnosis remains uncertain or requires specific expertise.
\end{itemize}

\section*{Treatment Options}
\begin{itemize}
  \item Treatment targets the underlying cause identified through diagnostic evaluation.
  \item Supportive care and symptom management are provided while awaiting definitive therapy.
  \item Medications, lifestyle modifications, and physical therapies may be prescribed as appropriate.
  \item Follow-up ensures treatment effectiveness and allows timely adjustments to the care plan.
\end{itemize}

\clearpage